\subsection{Fourier transform of tempered distributions}

Since every Schwartz function \(\varphi\) is integrable on \(\mathbf{R}^n\) (prop. 13 of IV, p. 213), it admits a Fourier transform \(\mathcal{F}(\varphi)\) (resp. a Fourier cotransform \(\overline{\mathcal{F}}(\varphi)\) ) which is identified with the continuous and bounded function on \(\mathbf{R}^n\) defined by

\[
y \mapsto \int_ {\mathbf {R} ^ {n}} \varphi (x) \exp (- 2 i \pi x \cdot y) d \mu (x)
\]

(resp. the function) \(y \mapsto \int_{\mathbf{R}^n} \varphi(x) \exp(2i\pi x \cdot y) d\mu(x)\) ).

Lemma 11. --- Let \(\varphi\in\mathcal{S}(\mathbf{R}^{n})\) . The function \(\mathcal{F}(\varphi)\) is infinitely differentiable on \(\mathbf{R}^{n}\) and we have

\[
\mathcal {F} (\mathrm{X} ^ {\alpha} \varphi) = (- 2 i \pi) ^ {- | \alpha |} \partial^ {\alpha} (\mathcal {F} (\varphi)),
\]

\[
\mathcal {F} (\partial^ {\alpha} \varphi) = (2 i \pi) ^ {| \alpha |} X ^ {\alpha} \mathcal {F} (\varphi)
\]

for all \(\alpha\) in \(\mathbf{N}^n\) .

We may assume that \(n \geqslant 1\) . Let \(\varphi \in \mathcal{S}(\mathbf{R}^{n})\) . The function defined by \((x, y) \mapsto \varphi(x) \exp(2i\pi x \cdot y)\) from \(\mathbf{R}^{n} \times \mathbf{R}^{n}\) into \(\mathbf{C}\) satisfies the hypotheses of Corollary 1 of IV, p. 198 for every integer k. The Fourier transform of \(\varphi\) is therefore infinitely differentiable and satisfies

\[
\partial^ {\alpha} (\mathcal {F} \varphi) (y) = (- 2 i \pi) ^ {| \alpha |} \int_ {\mathbf {R} ^ {n}} x ^ {\alpha} \varphi (x) \exp (- 2 i \pi x \cdot y) d \mu (x)
\]

for all \(y \in \mathbf{R}^{n}\) , which implies the first formula.

Let us prove the second formula. By induction on \(|\alpha|\) it suffices to do it when \(|\alpha| = 1\) , and one easily reduces to the case \(\alpha = (1,0,\ldots,0)\) . Let us write every \(x \in \mathbf{R}^n\) in the form \(x = (x_1,x')\) with \(x' \in \mathbf{R}^{n-1}\) , and denote by \(\mu_1\) (resp. \(\mu'\) ) the Lebesgue measure on \(\mathbf{R}^{n-1}\) (resp. \(\mathbf{R}\) ). By the Lebesgue--Fubini theorem (INT, V, p. 96, § 8, n° 4, th. 1), for every \(y = (y_1,y') \in \mathbf{R} \times \mathbf{R}^{n-1}\) , it follows that

\[
\begin{array}{l} \mathcal {F} (\partial_ {1} \varphi) (y) = \int_ {\mathbf {R} ^ {n - 1}} \exp (- 2 i \pi   x ^ {\prime} \cdot y ^ {\prime}) \\ \qquad \times \Big (\int_ {\mathbf {R}} (\partial_ {1} \varphi) (x _ {1}, x ^ {\prime}) \exp (- 2 i \pi   x _ {1} y _ {1}) d \mu_ {1} (x _ {1}) \Big) d \mu^ {\prime} (x ^ {\prime}). \end{array}
\]

For every compact interval \([a,b]\) in R and every \(x'\in R^{n-1}\) , we have

\[
\begin{array}{r l} & {\int_ {a} ^ {b} (\partial_ {1} \varphi) (x _ {1}, x ^ {\prime}) \exp (- 2 i \pi x _ {1} y _ {1}) d \mu_ {1} (x _ {1}) =} \\ & {\qquad \left[ \varphi (x _ {1}, x ^ {\prime}) \exp (- 2 i \pi x _ {1} y _ {1}) \right] _ {a} ^ {b}} \\ & {\qquad + 2 i \pi y _ {1} \int_ {a} ^ {b} \varphi (x _ {1}, x ^ {\prime}) \exp (- 2 i \pi x _ {1} y _ {1}) d \mu_ {1} (x _ {1})} \end{array}
\]

by integration by parts (FVR, II, p. 10, formula (10)). When \(a\) tends to \(-\infty\) and \(b\) tends to \(+\infty\) , the first term on the right-hand side converges to 0 since \(\varphi \in \mathcal{S}(\mathbf{R}^n)\) . The second term converges by Lebesgue's theorem (INT, IV, p. 137, § 3, n° 7, th. 6) to

\[
2 i \pi y _ {1} \int_ {\mathbf {R}} \varphi (x _ {1}, x ^ {\prime}) \exp (- 2 i \pi x _ {1} y _ {1}) d \mu_ {1} (x _ {1}),
\]

since the map \(x_{1} \mapsto \varphi(x_{1}, x')\) is integrable on R. Since \(x_{1} \mapsto \partial_{1}\varphi(x_{1}, x')\) is also integrable on R, it follows that

\[
\begin{array}{r l r} & & {\int_ {\mathbf {R}} \partial_ {1} \varphi (x _ {1}, x ^ {\prime}) \exp (- 2 i \pi x _ {1} y _ {1}) d \mu_ {1} (x _ {1})} \\ & & {= 2 i \pi y _ {1} \int_ {\mathbf {R}} \varphi (x _ {1}, x ^ {\prime}) \exp (- 2 i \pi x _ {1} y _ {1}) d \mu_ {1} (x _ {1})} \end{array}
\]

and finally, applying the Lebesgue--Fubini theorem again, we conclude that

\[
\mathcal {F} (\partial_ {1} \varphi) (y) = 2 i \pi y _ {1} \int_ {\mathbf {R} ^ {n}} \varphi (x) \exp (- 2 i \pi x \cdot y) d \mu (x),
\]

as desired.

PROPOSITION 18. --- The restriction to \(\mathcal{S}(\mathbf{R}^{n})\) of the Fourier transformation is an automorphism of topological vector spaces whose inverse is the restriction of the Fourier cotransformation.

Let \(\varphi\in\mathcal{S}(\mathbf{R}^{n})\) . By the preceding lemma, the Fourier transform of \(\varphi\) belongs to \(\mathcal{C}^{\infty}(\mathbf{R}^{n})\) . Moreover, for \(\alpha\in\mathbf{N}^{n}\) and \(\beta\in\mathbf{N}^{n}\) , we have

\[
\begin{array}{r} \mathrm{X} ^ {\beta} \partial^ {\alpha} (\mathcal {F} (\varphi)) = (- 2 i \pi) ^ {| \alpha |} \mathrm{X} ^ {\beta} \mathcal {F} (\mathrm{X} ^ {\alpha} \varphi) \\ = (- 1) ^ {| \alpha |} (2 i \pi) ^ {| \alpha | - | \beta |} \mathcal {F} (\partial^ {\beta} (\mathrm{X} ^ {\alpha} \varphi)). \end{array}
\]

In particular, the function \(\mathrm{X}^{\beta}\partial^{\alpha}(\mathcal{F}(\varphi))\) is bounded. Since \(\alpha\) and \(\beta\) are arbitrary in \(\mathbf{N}^{n}\) , this means that \(\mathcal{F}(\varphi)\) belongs to \(\mathcal{S}(\mathbf{R}^{n})\) . Moreover, this calculation implies

\[
q _ {\alpha , \beta} (\mathcal {F} (\varphi)) \leqslant (2 \pi) ^ {| \alpha | - | \beta |} \| \partial^ {\beta} (X ^ {\alpha} \varphi) \| _ {1}
\]

for \((\alpha, \beta) \in \mathbf{N}^n \times \mathbf{N}^n\) and \(\varphi \in \mathcal{S}(\mathbf{R}^n)\) .

Since the inclusion of the space \(\mathcal{S}(\mathbf{R}^{n})\) in \(\mathrm{L}^{1}(\mathbf{R}^{n})\) is continuous (prop. 13 of IV, p. 213), the map \(q_{\alpha,\beta} \circ \mathcal{F}\) of \(\mathcal{S}(\mathbf{R}^{n})\) in \(\mathbf{R}\) is continuous. It follows that the Fourier transformation is continuous from the space \(\mathcal{S}(\mathbf{R}^{n})\) into itself (cf. EVT, II, p. 7, prop. 5, c)). One verifies similarly that the Fourier cotransformation is continuous from the space \(\mathcal{S}(\mathbf{R}^{n})\) into itself. By the Fourier inversion formula (theorem 3 of II, p. 222), the Fourier transformation and the Fourier cotransformation are mutually inverse isomorphisms.

DEFINITION 6. --- The Fourier transformation (resp. Fourier cotransformation) on \(\mathcal{S}'(\mathbf{R}^n)\) is called the transpose of the Fourier transformation on \(\mathcal{S}(\mathbf{R}^n)\) (resp. of the Fourier cotransformation).

We still denote by \(\mathcal{F}\) (resp. \(\overline{\mathcal{F}}\) ) the Fourier transformation (resp. the Fourier cotransformation) on \(\mathcal{S}'(\mathbf{R}^n)\) . The Fourier transformation on \(\mathcal{S}'(\mathbf{R}^n)\) is therefore an automorphism of topological vector spaces whose inverse is the Fourier cotransformation. For every \(f \in \mathcal{S}'(\mathbf{R}^n)\) , the tempered distribution \(\mathcal{F}(f)\) (resp. \(\overline{\mathcal{F}}(f) )\) is defined by the formula

\[
\langle \mathcal {F} (f), \varphi \rangle = \langle f, \mathcal {F} (\varphi) \rangle \quad (\text { resp. } \langle \overline {{\mathcal {F}}} (f), \varphi \rangle = \langle f, \overline {{\mathcal {F}}} (\varphi) \rangle)
\]

for all \(\varphi\in\mathcal{S}(\mathbf{R}^{n})\) .

PROPOSITION 19. --- Let f be a tempered distribution associated with a bounded measure \(\nu \in \mathcal{M}^{1}(\mathbf{R}^{n})\) (resp. with \(g \in \mathrm{L}^{2}(\mathbf{R}^{n})\) ). The Fourier transformation of f in \(\mathcal{S}'(\mathbf{R}^{n})\) is the tempered distribution associated with the Fourier transformation of the measure \(\nu\) (resp. with the Fourier transformation of g).

Let \(\nu\) be a bounded measure on \(\mathbf{R}^{n}\) and let f be the tempered distribution associated with \(\nu\) . The Fourier transformation \(\mathcal{F}(\nu)\) is a continuous and bounded function on \(\mathbf{R}^{n}\) (prop. 3 of II, p. 207). The tempered distribution associated with this function satisfies

\[
\langle \mathcal {F} (\nu), \varphi \rangle = \int_ {\mathbf {R} ^ {n}} \mathcal {F} (\nu) \varphi d \mu = \int_ {\mathbf {R} ^ {n}} \mathcal {F} (\varphi) d \nu = \langle f, \mathcal {F} (\varphi) \rangle = \langle \mathcal {F} (f), \varphi \rangle
\]

for all \(\varphi\in\mathcal{S}(\mathbf{R}^{n})\) , where the second equality is prop. 13 of II, p. 221, which is applicable here since the measure \(\varphi\cdot\mu\) is bounded. The tempered distribution associated with \(\mathcal{F}(\nu)\) is therefore \(\mathcal{F}(f)\) .

When \(f\) is the tempered distribution associated with \(g \in \mathrm{L}^2(\mathbf{R}^n)\) , we follow a completely analogous procedure, using formula (29) of II, p. 221.

A similar statement also holds for the Fourier cotransformation.

Remark. --- The elementary formulas concerning the Fourier transformation of measures remain valid for the Fourier transformation of tempered distributions. For example, for \(f \in \mathcal{S}'(\mathbf{R}^{n})\) and \(\alpha \in \mathbf{N}^{n}\) , we have

\[
\begin{array}{l} \mathcal {F} (\partial^ {\alpha} f) = (2 i \pi) ^ {| \alpha |} X ^ {\alpha} \mathcal {F} (f) \\ \mathcal {F} (X ^ {\alpha} f) = (- 2 i \pi) ^ {- | \alpha |} \partial^ {\alpha} (\mathcal {F} (f)) \end{array}
\]

by Lemma 11,

